% preamble.tex — paquetes y macros de notación
% Apunte LyC (Reynolds, Theories of Programming Languages, 2009)

% --- Idioma y codificación (necesarios para acentos del texto en español) ---
\usepackage[utf8]{inputenc}
\usepackage[T1]{fontenc}
\usepackage[spanish]{babel}

% --- Matemática ---
\usepackage{amsmath}
\usepackage{amssymb}
\usepackage{amsthm}
\usepackage{mathtools}
\usepackage{stmaryrd}   % \llbracket \rrbracket para la semántica

% --- Árboles ---
\usepackage{bussproofs} % árboles de derivación
\usepackage{forest}     % árboles de sintaxis

% --- Enlaces ---
\usepackage{hyperref}
\hypersetup{
  plainpages=false,   % evita destinos PDF duplicados (page.1)
  colorlinks=true,
  linkcolor=blue!50!black,
  citecolor=blue!50!black,
  urlcolor=blue!50!black,
}

% --- Entornos de teoremas (en español) ---
\theoremstyle{plain}
\newtheorem{theorem}{Teorema}[chapter]
\newtheorem{proposition}[theorem]{Proposición}

\theoremstyle{definition}
\newtheorem{definition}[theorem]{Definición}
\newtheorem{example}[theorem]{Ejemplo}

% --- Notación ---
\newcommand{\sem}[1]{\llbracket #1 \rrbracket}
\newcommand{\semi}[1]{\sem{#1}_{\mathsf{intexp}}}
\newcommand{\semb}[1]{\sem{#1}_{\mathsf{boolexp}}}
\newcommand{\semc}[1]{\sem{#1}_{\mathsf{comm}}}
